\documentclass[11pt]{article}
\usepackage[margin=1in]{geometry}
\usepackage{amsmath,amssymb,amsthm,mathtools}
\usepackage{graphicx}
\usepackage[hidelinks]{hyperref}
\usepackage{xcolor}
\usepackage{microtype}
\newtheorem{theorem}{Theorem}
\newtheorem{proposition}{Proposition}
\newtheorem{corollary}{Corollary}
\theoremstyle{definition}
\newtheorem{definition}{Definition}
\newcommand{\BHF}{\operatorname{BHF}}
\newcommand{\UMD}{\mathrm{UMD}}
\title{A Coherent-Hilbert Subcase of the UMD Openness Conjecture\\
and Its Endpoint Bilinear-Hilbert Consequence}
\author{Candidate partial result for arXiv:1909.06416}
\date{13 August 2026}

\begin{document}
\maketitle

\begin{center}
\fcolorbox{black}{gray!8}{\parbox{0.91\linewidth}{
\textbf{Status.} Candidate partial; likely valid, subject to human review.
The argument proves Conjecture~8.9 for a coherent two-stage interpolation
subclass.  It does \emph{not} synchronize arbitrary Hilbert endpoints and does
not claim the full conjecture.}}
\end{center}

\section{Source question}

The source is A.~Amenta and G.~Uraltsev, \emph{The bilinear Hilbert transform
in UMD spaces} \cite{amenta-uraltsev}.  The target is Conjecture~8.9 and
Remark~8.10 on PDF page~67 of arXiv:1909.06416.  For
$r\in[2,\infty)$, a Banach space $X$ is called $r$-intermediate UMD when
there are a compatible couple $(Y,H)$, with $Y$ UMD and $H$ Hilbert, such
that
\[
  X\simeq [Y,H]_{2/r}
\]
for the complex interpolation functor.

\begin{figure}[ht]
  \centering
  \includegraphics[width=\linewidth]{figures/open_problem_crop.png}
  \caption{Full-width raster crop of source PDF page~67, containing the
  endpoint question, the complete statement of Conjecture~8.9, and its
  immediate BHT consequence.}
\end{figure}

\paragraph{Transcription.}
Conjecture~8.9 asks: if $X$ is $r$-intermediate UMD for some
$r\in(2,\infty)$, must there exist $\widetilde r<r$ such that $X$ is
$\widetilde r$-intermediate UMD?  The source observes that a positive answer
would treat its remaining endpoint regime
\begin{equation}\label{eq:endpoint}
  \sum_{i=1}^{3}\frac1{r_i}=1.
\end{equation}
Remark~8.10 sketches a reiteration proof but correctly warns that the Hilbert
spaces in two successive representations need not be the same compatible
endpoint.

\section{Idea of the proof}

The source's warning identifies the whole issue.  If the two stages really do
use one Hilbert endpoint $H$, they are not merely abstractly similar: they lie
on one complex interpolation scale.  Reiteration then combines their two
positive Hilbert weights.  The combined interpolation parameter is strictly
larger than $2/r$, hence its corresponding intermediate exponent is strictly
smaller than $r$.

For an endpoint BHT triple, improving a single $r_j>2$ makes the reciprocal
sum in \eqref{eq:endpoint} strictly greater than one.  This opens a genuine
region of H\"older exponents to which the source's Theorem~1.1 applies.

\section{The coherent-Hilbert theorem}

\begin{definition}
Let $r>2$ and put $\theta=2/r$.  An $r$-intermediate representation of $X$
is \emph{coherently twice intermediate} if there are one compatible couple
$(Z,H)$, where $Z$ is UMD and $H$ is Hilbert, and an $\alpha\in(0,1)$ such
that
\begin{equation}\label{eq:coherent}
  X\simeq \bigl[[Z,H]_{\alpha},H\bigr]_{\theta}.
\end{equation}
Thus the left endpoint in the outer representation is itself obtained by
moving toward the \emph{same} Hilbert endpoint.
\end{definition}

\begin{theorem}[strict improvement for a coherent endpoint]
Suppose \eqref{eq:coherent} holds.  Set
\begin{equation}\label{eq:beta}
  \beta=\theta+\alpha(1-\theta),
  \qquad \widetilde r=\frac{2}{\beta}.
\end{equation}
Then
\[
  2<\widetilde r<r
  \qquad\text{and}\qquad
  X\simeq[Z,H]_{2/\widetilde r}.
\]
In particular, $X$ satisfies the conclusion of the UMD openness conjecture at
$r$.
\end{theorem}

\begin{proof}
The complex reiteration theorem, in its endpoint form, gives
\[
 \bigl[[Z,H]_{\alpha},H\bigr]_{\theta}
   =[Z,H]_{(1-\theta)\alpha+\theta}
   =[Z,H]_{\beta};
\]
see Bergh--L\"ofstr\"om, Theorem~4.6.1 \cite{bergh-lofstrom}, which is also
the theorem invoked in source Remark~8.10.  Since
$0<\alpha<1$ and $0<\theta<1$,
\[
  \theta<\theta+\alpha(1-\theta)<1.
\]
Consequently $2<2/\beta<2/\theta=r$.  Since $Z$ is UMD and $H$ is
Hilbert, the last interpolation identity is exactly a
$\widetilde r$-intermediate UMD representation.
\end{proof}

For context, the easy direction of the exponent geometry is automatic.

\begin{proposition}[upward-ray property]
If $X$ is $r_0$-intermediate UMD, then it is $r$-intermediate UMD for every
$r>r_0$.
\end{proposition}

\begin{proof}
Write $X\simeq[Y,H]_{\theta_0}$ with $\theta_0=2/r_0$, and let
$\theta=2/r<\theta_0$.  Put
\[
  \delta=\frac{\theta_0-\theta}{1-\theta}\in(0,1),
  \qquad W=[Y,H]_{\delta}.
\]
The space $W$ is UMD because the UMD property is stable under complex
interpolation.  Reiteration gives
$[W,H]_{\theta}=[Y,H]_{(1-\theta)\delta+\theta}
=[Y,H]_{\theta_0}\simeq X$.
\end{proof}

This proposition shows that Conjecture~8.9 concerns whether the left endpoint
of this upward ray is ever attained above $2$.  The coherent theorem excludes
attainment for the stated subclass.

\section{Endpoint bilinear Hilbert transform consequence}

We use the notation of source Theorem~1.1.  Let $X_i$ be
$r_i$-intermediate UMD, let
$\Pi:X_1\times X_2\times X_3\to\mathbb C$ be a bounded trilinear form, and
let $\BHF_{\Pi}$ denote the associated trilinear bilinear-Hilbert form.

\begin{corollary}[nonvacuous endpoint subcase]
Assume \eqref{eq:endpoint}.  Suppose some index $j$ with $r_j>2$ has a
coherently twice intermediate representation, and let $\widetilde r_j$ be
given by \eqref{eq:beta}.  Define
\[
  s_j=\widetilde r_j,
  \qquad s_i=r_i\quad(i\ne j).
\]
Then $\sum_i s_i^{-1}>1$, and source Theorem~1.1 gives
\[
 |\BHF_{\Pi}(f_1,f_2,f_3)|
 \lesssim \prod_{i=1}^{3}\|f_i\|_{L^{p_i}(\mathbb R;X_i)}
\]
for every H\"older triple $p_i\in(1,\infty)$ satisfying
\begin{equation}\label{eq:source-exponents}
 \sum_{i=1}^{3}
 \frac{1}{\min(p_i,s_i)'(s_i-1)}>1.
\end{equation}
Moreover, the set of such H\"older triples is nonempty.
\end{corollary}

\begin{proof}
Only the $j$th exponent changes, so
\[
 \sum_i\frac1{s_i}
 =\frac1{\widetilde r_j}+\sum_{i\ne j}\frac1{r_i}
 >\sum_i\frac1{r_i}=1.
\]
The theorem above makes $X_j$ $s_j$-intermediate, while the other spaces
retain their original representations.  Source Theorem~1.1 therefore proves
the asserted bound whenever \eqref{eq:source-exponents} holds.

To see that this condition is not vacuous, choose numbers
$a_i\in(0,1/s_i]$ with $\sum_i a_i=1$; such a choice exists because
$\sum_i1/s_i>1$.  Put $p_i=1/a_i$.  Then $(p_i)$ is a finite H\"older
triple and $p_i\ge s_i$, so
\[
 \frac{1}{\min(p_i,s_i)'(s_i-1)}
 =\frac{1}{s_i'(s_i-1)}=\frac1{s_i}.
\]
Summation gives \eqref{eq:source-exponents}.
\end{proof}

\section{Stress tests and limitations}

\begin{itemize}
\item The strict inequalities in \eqref{eq:beta} use both $r>2$ and
      $0<\alpha<1$; no endpoint convention is hidden.
\item The endpoint BHT corollary uses the improved tuple $(s_i)$ in the
      source exponent condition.  With the original endpoint tuple, every
      summand in that condition is at most $1/r_i$, so its admissible region
      would be empty.
\item Scalar $L^p$ and Schatten interpolation scales provide consistency
      checks because their standard complex interpolation identities use one
      fixed $L^2$ or $S^2$ endpoint.  Bounds for those classical examples are
      already known and are not claimed as novel here.
\item The proof does not turn two abstract representations
      $X\simeq[Y_1,H_1]_{2/r}$ and
      $Y_1\simeq[Y_2,H_2]_{2/s}$ into a coherent one.  An isomorphism
      $H_1\simeq H_2$ need not extend to an isomorphism of compatible couples.
      That is the unresolved core of Conjecture~8.9.
\item Eight focused upgrade attempts were made: arbitrary endpoint
      isomorphisms, Hilbert stabilization, graph/pushout constructions,
      subquotient interpolation, type/cotype inversion, lattice factorization,
      endpoint limiting, and reversal of the tensor-BHT counterexample.  Each
      either reproduces the compatibility gap or fails an exact hypothesis;
      details are in the accompanying attempt note.
\end{itemize}

\section{Novelty and review recommendation}

Run indexes and bounded searches by arXiv id, exact conjecture name, source
title, and the phrases ``UMD openness'' and ``intermediate UMD'' found no later
claimed solution.  This packet should nevertheless be viewed as a rigorous
structural completion of the source's own Remark~8.10, not as a claim that the
full conjecture has been solved.  Novelty confidence is therefore low; the
mathematical value is the precise sufficient hypothesis and the explicit
nonvacuous endpoint-BHT corollary.

Human review should check (i) that the intended compatible-couple
identifications in \eqref{eq:coherent} are exact enough for endpoint
reiteration, and (ii) that the exponent region in the corollary is stated with
the improved tuple.  Once those points are accepted, the proof is a direct
application of the standard reiteration theorem.

\begin{thebibliography}{9}
\bibitem{amenta-uraltsev}
A. Amenta and G. Uraltsev,
\emph{The bilinear Hilbert transform in UMD spaces},
Math. Ann. 378 (2020), 1129--1221; arXiv:1909.06416.

\bibitem{bergh-lofstrom}
J. Bergh and J. L\"ofstr\"om,
\emph{Interpolation Spaces: An Introduction},
Grundlehren der mathematischen Wissenschaften 223, Springer, 1976.
\end{thebibliography}

\end{document}
